Refrigerant charge determination, pull-down, the performance test at
\(25^\circ\mathrm{C}\), and the three consumer-usage tests are required
for every product.
The temperature and humidity of each test are given in
Table~\ref{tab:test_attributes}.
Ambient humidity means that the test needs no humidity adjustment and
can be run in any chamber.
An entry of 85\% RH means that the test must be run at 85\% relative
humidity, so it can be assigned only to a chamber with humidity control.
Every test uses one sample, except pull-down, which is repeated on each
sample.
For refrigerant charge determination and pull-down, the temperature is
the upper ambient limit of the product's climatic class.
Products targeted for tropical, subtropical, and normal/subnormal
climates are tested at \(43^\circ\mathrm{C}\), \(38^\circ\mathrm{C}\),
and \(32^\circ\mathrm{C}\), respectively.
In addition to the performance test at \(25^\circ\mathrm{C}\), the
climatic class determines one low-temperature and one high-temperature
performance test, as specified in Table~\ref{tab:test_attributes}.
An additional performance test at \(16^\circ\mathrm{C}\) is required for
products targeting tropical, subtropical, and normal climates, whereas
the corresponding test for products targeting subnormal climates is
performed at \(10^\circ\mathrm{C}\).
High-temperature performance tests at \(43^\circ\mathrm{C}\),
\(38^\circ\mathrm{C}\), and \(32^\circ\mathrm{C}\) are required for
products targeting tropical, subtropical, and normal or subnormal
climates, respectively.
The applicable energy-efficiency tests are determined by the
energy-labeling regulation of the product's target market.
That regulation specifies exactly one of three mutually exclusive
options: the pair of tests at \(16^\circ\mathrm{C}\) and
\(32^\circ\mathrm{C}\), a single test at \(25^\circ\mathrm{C}\), or a
single test at \(32^\circ\mathrm{C}\).
The three energy-efficiency rows in Table~\ref{tab:test_attributes} are
therefore alternatives, not a mandatory set.
Freezing-capacity and temperature-rise tests are optional and are
included when requested by the product engineer.
Thus, each product generally has a different test portfolio even though
all products follow the same stage-based testing structure.
